\documentclass[11pt]{article}
\usepackage{amsmath,amssymb}
\usepackage[margin=1in]{geometry}
\usepackage{booktabs}
\usepackage{tabularx}
\usepackage{hyperref}
\setlength{\emergencystretch}{2em}

\title{General Open-World Scientific Discovery Superiority:\\
Discover More, Preserve Search Optionality, and Know When the Science Is Closed}
\author{ORION Paper II Ultimate-Successor Working Manuscript}
\date{August 2026}

\begin{document}
\maketitle

\begin{abstract}
Deep-research systems now combine broad query generation, iterative retrieval, evidence processing, coverage estimation, contradiction handling, state management, and verification-aware stopping. ORION Paper II therefore targets a claim above any one of those mechanisms: across heterogeneous scientific domains, source ecosystems, route geometries, and changing questions, can ORION outperform a donor-complete research controller at discovering decision-relevant evidence, preserving future search optionality, synthesizing conflicting evidence, and issuing calibrated task-global closure under matched total resources? Existing P2 false-closure, missed-route, retrieved-but-unprocessed, and provider-invalid outcomes remain immutable causal evidence. They are not converted into wins and are not used to shrink the target. Modern evidence-aware termination, claim-evidence coverage, retrieval-novelty control, high-fan-out exploration, coverage estimation, contradiction-aware synthesis, caching, and query-conditioned state compilation are granted to the strongest comparator. P2-U asks whether ORION can govern these donor mechanisms as scientific obligations, allocating effort across route generation, retrieval, processing/state preservation, and closure while failing closed on unresolved authority. The intended terminal is replicated \textbf{GENERAL\_OPEN\_WORLD\_DISCOVERY\_SUPERIORITY}, including at least one failure-derived exploration, state, or closure mechanic that transfers beyond its originating regime.
\end{abstract}

\section{Largest scientific question}
\begin{quote}
At matched model, tool, information, and total resource budgets, can ORION become a generally superior open-world scientific discovery process by finding more of the evidence that matters, retaining the ability to discover later-relevant evidence, and closing only when the task-level scientific obligations are actually discharged?
\end{quote}

The comparison frontier is demanding. DeepResearch Bench evaluates research-report quality and retrieval/citation behaviour on expert-authored tasks \cite{deepresearchbench2026}. LiveDRBench characterizes deep research through high-fan-out concept exploration and reports large query-coverage gaps across current systems \cite{livedrbench2026}. Enterprise Deep Research uses coverage-driven objectives and explicit evidence sufficiency criteria to reduce premature stopping \cite{dontstopearly2026}. HALT treats stopping as claim-evidence coverage rather than generator confidence \cite{halt2026}. RAAC uses retrieval novelty and information coverage to adapt deep-research trajectories and termination \cite{raac2026}. SAGE shows that reasoning-intensive scientific retrieval remains difficult and that retrieval rankings can change sharply with scientific task structure \cite{sage2026}. P2-U treats all of these as donor pressure, not as residual novelty to evade.

\section{Upward-generalization principle}
Any donor mechanism that solves a P2-U subclass is added to the comparator. ORION earns the broader result only through a combination of:
\begin{enumerate}
    \item higher decision-relevant discovery under matched total resources;
    \item lower false task-global closure at equal or higher scientific utility;
    \item better adaptation to heterogeneous route, evidence, and contradiction regimes;
    \item preserved optionality under future question and evidence-horizon changes frozen before current-state construction;
    \item cross-domain transfer of the governing search-and-closure policy; and
    \item at least one failure-derived exploration, state, or closure mechanic that expands protected reach beyond the pre-freeze donor-complete basis.
\end{enumerate}
Coverage estimation, value-of-information, adaptive stopping, query novelty, contradiction handling, and state compression are not claimed individually as ORION inventions.

\section{Immutable adverse evidence}
The predecessor P2 external probes include non-positive and provider-invalid outcomes. The successor preserves the distinctions among:
\begin{enumerate}
    \item a material route never generated;
    \item evidence never retrieved;
    \item evidence retrieved but excluded incorrectly;
    \item evidence included but not processed into decision-usable state;
    \item state compression that removes a future route cue;
    \item provider or tool output invalid or unavailable;
    \item local route saturation promoted incorrectly to task-global closure.
\end{enumerate}
A provider-invalid episode remains \textsc{CannotCheck}. It is never scored as a comparator loss. Historical false closure remains false closure even if a successor later fixes the responsible mechanism.

\section{Scientific-obligation task model}
Each protected task contains hidden scientific obligations $O=\{o_1,\ldots,o_m\}$ and an evidence universe partitioned into candidate route families $R$. Candidate systems do not receive the gold obligation set, route set, or materiality labels. At time $t$ the system maintains
\[
S_t=(Q_t,E_t,P_t,U_t,C_t),
\]
where $Q_t$ is the active route/query frontier, $E_t$ acquired raw evidence, $P_t$ processed evidence state, $U_t$ unresolved obligations and preserved optionality cues, and $C_t$ candidate closure status.

The protected evaluator distinguishes four task terminals:
\begin{itemize}
    \item \textsc{Closed}: every load-bearing obligation is discharged under the frozen criterion;
    \item \textsc{Open}: at least one identified material obligation remains actionable;
    \item \textsc{Unresolved}: available evidence cannot establish closure or a specific next material route within the frozen budget;
    \item \textsc{CannotCheck}: required provider, evidence custody, or evaluator authority is unavailable.
\end{itemize}
Search stagnation, confidence, citation count, context length, or duplicate retrieval is never itself closure authority.

\section{Failure-identifying benchmark families}
The protected benchmark is frozen before outcomes and contains surface-confusable pairs.
\begin{table}[h]
\centering
\small
\begin{tabularx}{\textwidth}{@{}lXX@{}}
\toprule
Family & Hidden condition & Failure exposed \\
\midrule
Missing route & decisive evidence is behind an ungenerated route & route generation/coverage \\
Retrieved but unprocessed & decisive source was fetched but not incorporated & processing state \\
Duplicate evidence & many citations map to one underlying evidence unit & false coverage inflation \\
Contradictory evidence & high-quality sources disagree materially & synthesis/closure calibration \\
Provider invalid & mandatory route cannot be verified & authority, not performance \\
Apparently saturated & repeated retrieval duplicates while hidden route remains & premature stopping \\
Optionality loss & compact state answers now but deletes future-route cue & state preservation \\
Underspecified/unanswerable & current question does not license a determinate answer & unresolved calibration \\
Clean closure & all obligations are truly discharged & excessive-abstention penalty \\
\bottomrule
\end{tabularx}
\caption{Stage-identifying P2-U families. Each positive-looking shortcut has a paired falsifier.}
\end{table}

Exact authored cases establish mechanism non-vacuity only. The headline study must include naturalistic scientific tasks across multiple source ecosystems, at least one systematic-review-like evidence set with independently curated relevance, and a block that jointly holds out domain and route geometry.

\section{Obligation-Directed Adaptive Exploration}
The first negative-to-positive successor family is \textbf{Obligation-Directed Adaptive Exploration (ODAE)}. ODAE is not claimed as invention of active search or value-of-information. It composes donor-owned search controls with explicit scientific-obligation semantics.

At each step, the controller may allocate charged budget to route generation, retrieval, evidence processing, contradiction resolution, state compilation/cache/recovery, or closure verification. Candidate actions are ranked by expected reduction in decision-relevant obligation uncertainty per total charged cost. A state-compression action is inadmissible when it erases a frozen optionality coordinate whose loss cannot be recovered within the remaining budget. Task closure is admissible only when every currently load-bearing obligation is independently discharged or is recorded as \textsc{Unresolved}/\textsc{CannotCheck} with the corresponding authority boundary.

The research residual is whether this cross-stage responsibility composition outperforms a donor-complete controller on fresh regimes and whether failures yield a new transferable obligation/route/state mechanic.

\section{Donor-complete baselines}
All systems receive equivalent model class, retrieval/tool access, candidate-visible information, and total resource budgets. The mandatory baseline set includes:
\begin{enumerate}
    \item fixed-budget search+RAG;
    \item iterative deep-research agent;
    \item high-fan-out exploration controller informed by LiveDRBench;
    \item EDR-style evidence-aware termination \cite{dontstopearly2026};
    \item HALT-style claim-evidence coverage stopping \cite{halt2026};
    \item RAAC-style novelty/coverage retrieval control \cite{raac2026};
    \item SAGE-informed scientific retrieval system \cite{sage2026};
    \item coverage-estimator/capture-recapture augmented search;
    \item contradiction-aware synthesis;
    \item query-conditioned state compilation with cache/recovery policy;
    \item donor-complete controller containing all preceding mechanisms and allowed to route among them;
    \item strongest runnable deep-research/scientific-search system available at protocol freeze;
    \item oracle obligation/route ceiling, reported only as an analysis ceiling.
\end{enumerate}
Unavailable official systems are labelled unavailable. Mechanism-matched reimplementations are labelled proxies and are not presented as official implementations.

\section{Primary hypotheses}
\paragraph{H1, generalized discovery superiority.}
Across the heterogeneous protected benchmark, ORION improves final scientific utility and decision-relevant evidence recall over the strongest runnable donor-complete comparator by prospectively frozen margins.

\paragraph{H2, closure superiority.}
ORION reduces false task-global closure while preserving prospectively frozen clean-closure non-inferiority or superiority.

\paragraph{H3, optionality superiority.}
Under future queries/routes frozen before current state construction, ORION preserves decision-relevant search optionality at better utility-cost trade-offs.

\paragraph{H4, cross-regime transfer.}
The H1--H3 direction survives held-out domains, source ecosystems, route geometries, contradiction structures, and query shifts under independent route accounting.

\paragraph{H5, mechanism expansion.}
At least one retained negative/tied family yields a new exploration, state, obligation, or closure mechanic that expands protected reach beyond the pre-freeze donor-complete basis and transfers to fresh tasks.

\section{Outcomes and statistical plan}
The independent unit is the predeclared research task or, where dependence requires it, a prospectively declared task cluster. Repeated model samples, trajectories, retrieved documents, and reminted synthetic variants are technical repeats and do not inflate $n$.

Primary outcomes are reported separately:
\begin{enumerate}
    \item decision-relevant evidence recall and precision;
    \item protected final scientific utility;
    \item false closure rate;
    \item true clean-closure coverage;
    \item calibrated \textsc{Unresolved} and \textsc{CannotCheck};
    \item missed-obligation count;
    \item optionality loss under frozen future queries;
    \item cross-domain/source/route transfer;
    \item tokens, retrieval calls, state-construction/storage cost, evaluator calls, wall time, and time-to-valid-closure.
\end{enumerate}
Before protected access, freeze superiority and non-inferiority margins, task/cluster hierarchy, family weighting, interval method, multiplicity policy, missing/provider-invalid handling, and catastrophic false-closure rules. Effect sizes and uncertainty are primary; significance-only language is insufficient.

\section{Framework correspondence}
The executable framework is canonical for implemented mechanics. P2-U is synchronized with framework version `0.3.9-shadow` and K/W/M. Current substrate includes \texttt{SEARCH.v1}, \texttt{SATURATE\_BOUNDED.v3}, \texttt{RetrievalBenchmark.rakl-v1}, \texttt{RouteFamilyHealth.rakl-v1}, \texttt{EpistemicContextCompiler.rakl-v1}, and \texttt{ExternalEvidenceManifest.v1}. Route-family health is telemetry only and grants no truth, abandonment, or scientific authority.

ODAE, naturalistic hidden-route gold, a learned cross-domain joint controller, and GENERAL\_OPEN\_WORLD\_DISCOVERY\_SUPERIORITY are prospective. They are not described as current runtime capabilities and are not added to the canonical mechanics registry by this manuscript.

\section{Negative-result recursion}
Every false closure, missed material source, harmful state compilation, and tied/negative family is preserved and assigned a primary responsibility from
\[
\{\text{route generation},\text{retrieval},\text{inclusion},\text{processing},\text{state optionality},\text{stopping},\text{provider},\text{global authority},\text{benchmark}\}.
\]
The programme then freezes a discriminator, searches materially distinct successor mechanisms and parent disciplines, and tests the proposed repair on fresh tasks. Outcome deletion, post-hoc endpoint switching, margin relaxation, baseline weakening, and counting provider failure as a competitor loss are forbidden.

ORION issue \#650 is the authority-bearing P2-U lane. The ChatGPT-led negative-recovery tranche uses ODAE as the first successor family. A subsequent failure opens another materially distinct mechanism family rather than lowering the final target by default.

\section{Claim ladder and current status}
\begin{enumerate}
    \item \textbf{Protocol frozen}: no performance claim.
    \item \textbf{Protected open-world superiority}: H1 and H2 pass with clean-closure safety.
    \item \textbf{Cross-regime discovery superiority}: H1--H4 pass across held-out domains/source/route regimes.
    \item \textbf{General open-world discovery superiority}: cross-regime superiority plus H5 demonstrates failure-driven expansion of the search/closure mechanism basis.
\end{enumerate}
This TeX is a prospectively refined successor. It does not rewrite predecessor results and does not assert the headline terminal before protected external evidence exists.

\begin{thebibliography}{12}
\bibitem{deepresearchbench2026}
Y. Yang et al., ``DeepResearch Bench: A Comprehensive Benchmark for Deep Research Agents,'' ICLR 2026.

\bibitem{livedrbench2026}
A. Java, A. Khandelwal, S. Midigeshi, A. Halfaker, A. J. Deshpande, N. Goyal, A. Gupta, N. Natarajan, and A. Sharma, ``Characterizing Deep Research: A Benchmark and Formal Definition,'' ICLR 2026.

\bibitem{dontstopearly2026}
P. K. Choubey et al., ``Don't Stop Early: Scalable Enterprise Deep Research with Controlled Information Flow and Evidence-Aware Termination,'' Proceedings of ACL 2026, Industry Track, pp. 1699--1713, doi:10.18653/v1/2026.acl-industry.116.

\bibitem{halt2026}
D. Roh and D. Han, ``HALT: Verification-Aware Stopping for Retrieval-Augmented Search Agents,'' arXiv:2608.02009, 2026.

\bibitem{raac2026}
H. Soudani, E. Lingg, F. Hasibi, and N. Rekabsaz, ``When Deep Research Agents Stagnate: Enhancing Reasoning with Retrieval-Aware Agent Control,'' arXiv:2608.15191, 2026.

\bibitem{sage2026}
T. Hu, Y. Zhao, C. Zhang, A. Cohan, and C. Zhao, ``SAGE: Benchmarking and Improving Retrieval for Deep Research Agents,'' arXiv:2602.05975, 2026.
\end{thebibliography}

\end{document}
